\documentclass[12pt, a4paper]{academics}

\academicsCoverLabel{解题报告}
\academicsCourse{算法设计与分析}{Design and Analysis of Algorithms}
\academicsReport{二分与分治}{2026 年 6 月 9 日}
\academicsArthur{华博文}{19240212}

\begin{document}

\makeacademicscover

\section{引言}

二分查找与分治法是算法设计中两个根基性的方法论。二分查找依托单调序列的结构性质，通过每次比较将搜索空间折半缩小，以对数代价完成定位；分治法则将原问题递归地分解为规模更小的同构子问题，独立求解后再将子问题的解合并为原问题的解。二者看似分属不同的算法范式，实则共享同一套思想内核——通过不断缩减问题规模来压缩时间复杂度。

本报告以四道难度递进、算法贯通的题目为载体，展示二分与分治从基础入门到高阶应用的完整图景。题一从有序数组上的二分查找出发，建立对单调性与 $\mathit{lower\_bound}$ 语义的基本理解，是整个报告的方法论起点。题二将二分的应用场景从「查找」拓展至「判定」——在木棍切割问题中，通过二分答案将难以直接求解的优化问题转化为一系列可高效判定的可行性问题，展现了二分作为一种元算法的通用性。题三正式进入分治领域，以归并排序为框架统计序列中的逆序对，重点展示分治过程中跨区间信息的合并技巧。题四则将分治思想推向多维空间——CDQ 分治在对三维偏序的处理中，将第一维交予初始排序、第二维交予归并中的双指针遍历、第三维交予树状数组，三种机制各司其职、层层降维，是分治思想与数据结构协同作战的典范。

报告中的所有代码均使用 C++23 标准编写，在 Apple clang 21.0.0 编译器下通过测试。

\section{题一：二分查找}

\subsection{问题重述}

给定一个长度为 $n$ 的严格非降整数序列 $a_1 \le a_2 \le \dots \le a_n$，以及 $k$ 个查询值。对于每个查询值，需要输出它在序列中第一次出现的位置（位置从 $1$ 开始计数）；若该值不存在于序列中，则输出 $-1$。数据范围为 $1 \le n, k \le 10^5$，序列元素和查询值的绝对值均不超过 $10^9$。

\subsection{单调性驱动的折半搜索}

二分查找之所以高效，根源在于序列的单调性赋予了其一种全局偏序结构：对任意位置 $i$，其左侧的所有元素均不大于 $a_i$，右侧的所有元素均不小于 $a_i$。这一性质意味着，一次比较便足以安全地淘汰一半的搜索空间。具体而言，任取当前搜索区间的中间位置 $m = \lfloor (l + r) / 2 \rfloor$，将 $a_m$ 与目标值比较：若 $a_m$ 严格小于目标值，则目标值不可能出现在 $[l, m]$ 内——因为该区间内的元素均不大于 $a_m$，从而严格小于目标值，搜索区间缩小为 $[m+1, r]$；若 $a_m$ 大于或等于目标值，则首次出现位置必然不会落在 $[m+1, r]$ 内——因为 $m$ 及其左侧才可能包含更早的出现位置，搜索区间缩小为 $[l, m]$。每次迭代区间长度严格减半，至多进行 $\lceil \log_2 n \rceil$ 次比较即可完成定位。

上述过程正是 $\mathit{lower\_bound}$ 的经典语义：在有序序列中定位第一个大于或等于目标值的元素。C++23 标准库中的 $\mathtt{ranges::lower\_bound}$ 直接封装了这一算法。对于每次查询，调用 $\mathit{lower\_bound}$ 获取迭代器后，检查其是否抵达序列末尾以及指向的元素是否等于目标值：若两个条件均满足，则该位置即为答案（转为 1-based 索引输出）；否则目标值不存在，输出 $-1$。值得指出的是，二分查找在手动实现时存在一个常见陷阱——左闭右开区间 $[l, r)$ 与闭区间 $[l, r]$ 的开闭语义混淆极易导致死循环或边界遗漏。使用标准库实现不仅规避了此类边界错误，也使代码的意图表达更为清晰。

每次二分查找的时间复杂度为 $\mathrm{O}(\log n)$，$k$ 次查询总计 $\mathrm{O}(k \log n)$。在 $n = k = 10^5$ 的极限数据下，约 $1.7 \times 10^6$ 次比较操作，可在毫秒级别内完成。空间复杂度为 $\mathrm{O}(n)$，仅需存储输入数组。

\subsection{代码实现}

\inputminted{c++}{../src/p1_binary_search.cpp}

\section{题二：二分答案}

\subsection{问题重述}

现有 $n$ 根木棍，第 $i$ 根的长度为 $L_i$。需要将这些木棍切割成至少 $k$ 段，每段必须来自同一根木棍且长度为整数，求可能的最长段长。若无论如何切割都无法得到 $k$ 段，则输出 $0$。数据范围：$1 \le n \le 10^5$，$1 \le k \le 10^9$，$1 \le L_i \le 10^9$。

\subsection{可行性单调性与二分答案范式}

本题初看是一个优化问题——直接求解最大段长似乎缺乏抓手。若枚举所有可能的段长并逐一检验，段长的取值范围可达 $10^9$，显然不可接受。然而，注意到问题中蕴藏着一个关键的单调性质：若段长 $\ell$ 是可行的（即可以切出至少 $k$ 段），那么任意小于 $\ell$ 的段长也必定可行——段长越短，每根木棍所能切出的段数只增不减。换言之，可行性函数关于段长是单调的：存在一个临界值 $\ell^*$，使得所有 $\ell \le \ell^*$ 均可行，而所有 $\ell > \ell^*$ 均不可行。可行域恰好是 $[1, \ell^*]$ 这样一个连续前缀。

这一单调性正是二分答案范式的入场券。我们将原优化问题转化为一族判定问题：对于候选段长 $\mathit{mid}$，判断能否切出至少 $k$ 段。判定函数只需遍历全部 $n$ 根木棍，计算 $\sum_{i=1}^{n} \lfloor L_i / \mathit{mid} \rfloor$ 并与 $k$ 比较即可，单次判定耗时 $\mathrm{O}(n)$。二分的搜索范围为 $[1, \max L_i]$：段长最小为 $1$（整数约束），最大不可能超过任意一根木棍的长度（否则该木棍无法贡献任何段）。在二分过程中，若 $\mathit{mid}$ 可行，则记录答案并向右搜索更大的段长；若不可行，则向左收缩区间。最终记录的答案即为最大可行段长。若在整段二分过程中从未遇到可行值（即段长 $1$ 时仍不可行），输出 $0$。

二分答案的价值远远超出了本题的范畴。在算法竞赛与工程实践中，凡具备「答案单调、判定容易」特征的问题——最小化最大距离、最大化最小收益、分配与调度问题等——均可用这一范式求解。它体现了一种重要的算法思维：当直接求解优化问题遇到瓶颈时，退一步构造可高效判定的子问题，再利用单调性二分逼近答案，往往能够化繁为简。

判定函数每次遍历全部 $n$ 根木棍，耗时 $\mathrm{O}(n)$。二分搜索共进行 $\mathrm{O}(\log \max L_i)$ 次判定，由于 $\max L_i \le 10^9 < 2^{30}$，迭代次数不超过 $30$。总时间复杂度为 $\mathrm{O}(n \log \max L_i)$，在 $n = 10^5$ 时约 $3 \times 10^6$ 次运算。空间复杂度为 $\mathrm{O}(n)$，仅需存储木棍长度。

\subsection{代码实现}

\inputminted{c++}{../src/p2_cutting_sticks.cpp}

\section{题三：归并排序与逆序对}

\subsection{问题重述}

给定一个长度为 $n$ 的整数序列 $a_1, a_2, \dots, a_n$。逆序对是指满足 $i < j$ 且 $a_i > a_j$ 的有序对 $(i, j)$。需要求出序列中逆序对的总数。数据范围：$1 \le n \le 10^5$，序列元素的绝对值不超过 $10^9$。

\subsection{分治框架与归并中的跨越统计}

朴素做法是枚举所有有序对 $(i, j)$ 并逐一检验 $a_i > a_j$，时间复杂度为 $\mathrm{O}(n^2)$，在 $n = 10^5$ 时达到百亿级别，完全不可行。分治策略则能将复杂度压缩至 $\mathrm{O}(n \log n)$。

算法的核心在于将逆序对按位置关系划分为三类。将当前区间 $a[l..r)$ 从中点 $\mathit{mid}$ 一分为二，得到左半区间 $a[l..\mathit{mid})$ 和右半区间 $a[\mathit{mid}..r)$。原区间内的所有逆序对必然属于以下三类之一：两个元素均在左半、两个元素均在右半、以及一个位于左半而另一个位于右半的「跨越」逆序对。前两类由递归调用在左半和右半上分别求解，第三类则嵌入归并排序的合并过程中进行统计。

归并两个已排序的子数组时，双指针 $i$ 在左半滑动、$j$ 在右半滑动。当 $a[i] \le a[j]$ 时，$a[i]$ 与右半中 $j$ 之前的所有元素均不构成逆序对（因为那些元素均不小于 $a[i]$），直接将 $a[i]$ 放入合并结果；当 $a[i] > a[j]$ 时，由于左半已升序排列，$a[i]$ 及其后所有左半元素均严格大于 $a[j]$，共构成 $\mathit{mid} - i$ 个跨越逆序对，将此数量累加至逆序对总数，随后将 $a[j]$ 放入合并结果。合并完成后，当前区间变为有序，同时跨越逆序对也全部统计完毕。

这一设计的精妙之处在于排序与计数完全同步。归并排序本身需要 $\mathrm{O}(n \log n)$ 的时间，逆序对统计的额外开销仅为合并过程中的一个常数次加法，不改变渐近复杂度。这正是分治法「在合并步骤中收集跨越信息」这一经典范式的体现。需要注意，逆序对总数在最坏情况下可达 $n(n-1)/2 \approx 5 \times 10^9$，超出 32 位有符号整数的表示范围，因此必须使用 64 位长整型（$\mathtt{long\ long}$）存储计数结果。归并的递归深度为 $\mathrm{O}(\log n)$，每层合并所有元素恰好一次，总时间复杂度 $\mathrm{O}(n \log n)$；由于需要临时数组辅助合并，空间复杂度为 $\mathrm{O}(n)$。

\subsection{代码实现}

\inputminted{c++}{../src/p3_inversions.cpp}

\section{题四：CDQ 分治（三维偏序）}

\subsection{问题重述}

给定 $n$ 个三元组，第 $i$ 个三元组为 $(a_i, b_i, c_i)$。对于每个三元组 $i$，需要统计满足 $a_j \le a_i$ 且 $b_j \le b_i$ 且 $c_j \le c_i$ 的三元组 $j$ 的数量（包含 $i$ 自身）。输出每个三元组的偏序计数。数据范围：$1 \le n \le 10^5$，所有坐标均为不超过 $10^9$ 的正整数。

\subsection{CDQ 分治的逐维消解策略}

多维偏序统计是算法竞赛中的经典问题类别。若仅有一维偏序约束，排序后计算前缀和即可 $\mathrm{O}(n)$ 求解；若有两维，排序消去第一维后，对第二维使用树状数组或归并排序即可 $\mathrm{O}(n \log n)$ 求解。三维偏序则需要更精巧的处理——CDQ 分治，由陈丹琦在 2008 年国家集训队论文中提出，正是为此类问题量身打造的方法。

CDQ 分治的核心思路是「以时间换维度」：将第一维的排序作为分治的「时间轴」，递归地处理左右两个区间。初始时将所有三元组按 $a$ 升序排列（$a$ 相同时依次按 $b$、$c$ 升序）。在分治的递归结构中，每一步将当前区间从中点划分为左半和右半，递归保证了左半区间中所有三元组的 $a$ 均不大于右半区间中所有三元组的 $a$——第一维的偏序约束由此被递归结构天然吸收。第一维消去之后，问题在每个递归层上退化为一个二维偏序问题：对于右半的每个三元组，统计左半中满足 $b$ 和 $c$ 约束的元素个数。剩余两维的统计通过归并排序与树状数组协同完成：在归并左右区间（按 $b$ 排序）的过程中，用 BIT 维护已经遍历过的左半元素的 $c$ 坐标，从而实现高效的二维查询。

三个维度就这样被分配给三种机制，每种机制贡献一个对数因子：初始排序处理 $a$ 维，归并的递归深度处理 $b$ 维，BIT 的每次操作处理 $c$ 维，三者叠加达到 $\mathrm{O}(n \log^2 n)$ 的总复杂度。这一设计体现了算法工程中「各司其职、分而治之」的哲学——面对多维复合约束，不是试图寻找一种全能的通解，而是将各维度的负担分别交由最擅长处理该维度的工具承担。

\subsection{归并流程与树状数组的协同}

算法的具体流程可分为三个步骤。第一步，按 $(a, b, c)$ 的优先顺序对所有三元组升序排列，为分治提供第一维的顺序保证。第二步，对 $c$ 坐标进行离散化压缩：由于 $c$ 的取值范围可达 $10^9$，而 BIT 的大小必须与值域相当，因此对所有 $c$ 值排序去重后，将每个 $c_i$ 映射为其在有序去重序列中的排名（$1$-based），使 BIT 的大小压缩至 $\mathrm{O}(n)$。第三步，执行 CDQ 分治。

分治函数处理区间 $[\ell, r)$。若区间长度不超过 $1$ 则直接返回；否则取中点 $m$，递归处理 $[\ell, m)$ 与 $[m, r)$。递归返回后，$[\ell, m)$ 和 $[m, r)$ 各自已按 $b$ 有序。合并时使用双指针 $i$（遍历左半）和 $j$（遍历右半）：当左半元素的 $b$ 不大于右半元素的 $b$ 时，将该左半元素的 $c$ 值在 BIT 中加 $1$（单点修改），并推进 $i$；当右半元素的 $b$ 更小时，在 BIT 中查询不超过该右半元素 $c$ 值的前缀和——这个前缀和即为满足三维偏序（$a$ 由递归保证，$b$ 由归并顺序保证，$c$ 由 BIT 查询保证）的左半元素个数，将其累加至该右半元素的答案中，并推进 $j$。合并完成后，清除 BIT 中本轮左半所有元素的贡献（回溯时不影响其他分支），并将合并结果写回原数组，使当前区间按 $b$ 有序以供父层归并使用。

每个三元组的答案初始化为 $1$（计入自身），CDQ 过程中累加的所有贡献均来自三维均不大于它的左侧三元组，最终输出即为完整的偏序计数。CDQ 分治的递归深度为 $\mathrm{O}(\log n)$，每层合并时所有元素恰好被处理一次，每次处理涉及一次 BIT 操作（$\mathrm{O}(\log n)$），总时间复杂度为 $\mathrm{O}(n \log^2 n)$。在 $n = 10^5$ 时，$n \log^2 n \approx 2.8 \times 10^7$，可在几百毫秒内完成。空间复杂度为 $\mathrm{O}(n)$，主要开销来自存储三元组、答案数组、BIT 以及合并时的临时数组。

\subsection{代码实现}

\inputminted{c++}{../src/p4_cdq_3d.cpp}

\section{总结}

本报告以二分与分治为主线，选取四道难度递进、思想贯通的问题，系统展示了该方法论从基础到高阶的完整演进脉络。

题一的二分查找是整个线索的起点。其理论基础在于单调序列所赋予的偏序结构——仅凭一次比较即可安全地淘汰一半搜索空间。二分查找并非仅仅是标准库中一个函数的调用方式，更是一种通用的搜索策略：只要搜索空间具备单调性，二分就能将对数级别的加速嵌入到任何场景之中。

题二将二分从「查找」升维为「答案搜索」。在木棍切割问题中，最优段长本身并不以显式形式存在于输入数据中，它是搜索空间中的一个待确定数值。二分答案范式将原优化问题重构为一族判定问题——对每个候选答案判断其可行性，再依据单调性收缩搜索区间。这一思维转换是算法设计中的核心能力：跳出直接求解优化的思维定式，转而构造可高效判定的子问题，往往能实现复杂度的质变。

题三正式进入分治领域。逆序对问题天然契合分治的条件：原问题可分解为两个半区间上的同构子问题，而合并步骤仅需在归并排序的过程中稍作扩展，即可完成跨越逆序对的统计。本题揭示了分治法的一个核心价值——在合并阶段，两个已经各自解决的子问题之间往往存在着可以高效计算的结构性信息，这些信息正是连接子问题解与整体答案的关键桥梁。

题四的 CDQ 分治将分治思想推至多维偏序场景。三个维度分别分配给三种机制——排序处理第一维、归并的时序处理第二维、树状数组处理第三维——每种机制在各自负责的维度上提供对数级别的效率，三者协同实现了 $\mathrm{O}(n \log^2 n)$ 的总复杂度。这种「逐维降阶」的策略具有广泛的适用性：在更高维度或更复杂约束的问题中，CDQ 分治可以进一步与树套树、K-D 树等数据结构结合，灵活应对不同维度的偏序需求。

回顾四道题目的递进关系，可以提炼出一条贯穿始终的方法论线索：二分提供了单维度上的对数加速，分治提供了问题规模上的对数缩减，而 CDQ 则是二者的有机融合——在分治框架中，通过归并（实质上是一种有序合并的二分过程）和树状数组（其索引结构本身就构建于二分的原理之上）协同运作，将多维约束层层消解。从二分的「折半」到分治的「分解」，再到 CDQ 的「逐维降阶」，同一思想在不同复杂度层面的延伸与展开，正是算法设计中最具深度与美感的部分。

\end{document}
